\chapter{Lie Groups and Lie Algebras}

 [...]

From the Heine Borel theorem we know that a subset of $\mathbb{R}$ is compact if and only if it is closed (there are no limiting points that belong to the set) and bounded. We can use this to show that the set of all subgroups of $\mathbb{R}$ is uncountable.

\begin{example}[Compactness of the unitary group]
    The unitary group \(G = U(n)\) satisfies the following properties:
    \[
        \sum_{i} u^*_{ij} u_{ik} = \delta_{jk},
    \]
    and expanding the sum we understand that for \(i=j\) we have no entries greater than 1:
    \[
        \vert u_{ij} \vert \leq 1.
    \]
    Furthermore if we have some matrix group defined by a set of polynomial equations, then the set of solutions is closed. Therefore \(U(n)\) is a closed and bounded subset of \(\mathbb{R}^{n^2}\) and thus compact.
\end{example}

\begin{example}[Non compact Lie Groups]
    \(GL(n,\mathbb{R})\) or \(GL(n,\mathbb{C})\), \(SL(n,\mathbb{R})\), \(SO(1,n-1)\) are not compact, since they are not bounded.
\end{example}

\begin{definition}[Connectedness]
    A topological space \(G\) is \textit{connected} if for any two ponts \(g_1,\,g_2 \in G\) there exists a continuous path \(\gamma : [0,1] \to G\) such that \(\gamma(0) = g_1\) and \(\gamma(1) = g_2\) (it is not necessary to impose the limits on the domanin of \(\gamma\), but it is convenient to do so).
\end{definition}

This curve is a continuous map, thus any function cannot vary discontinuously along it. In particular, if we have a continuous homomorphism \(f : G \to H\) and \(G\) is connected, then the image of \(f\) must be connected as well. This is because the image of a continuous map is always connected.

If we start from a point where \(\det(f(g)) = 1\) for some \(g \in G\), then we can infere that \(\det(f(g)) = 1\) for all \(g \in \gamma \subset G\). This is because the determinant is a continuous function, and if it is unitary at one point, it must be unitary along the entire path connecting that point to any other point in \(G\).

    [... unitary and orthogonal groups ...]

\begin{definition}[Simply connected Lie groups]
    If a group \(G\) is connected and any closed curve in \(G\) can be continuously contracted (smoothly deformed) to a point, then \(G\) is called \textit{simply connected}.
\end{definition}

It is easier to think about topological spaces as being made of rubber, and the continuous deformation as being a process of stretching and bending the space without tearing it. In this way, a simply connected space is one that has no holes or loops that cannot be contracted to a point. As a visual example, the case of a thorus, a can connect any two points with a continuous path, but there are loops around the hole of the thorus that cannot be contracted to a point. On the other hand, a sphere is simply connected, because any loop on the sphere can be contracted to a point.

\begin{example}[Simply connectedness in Lie groups]
    The unitary group \(U(n)\) is not simply connected, since it contains \(U(1) \subset U(n)\) which is topologically a circle, and thus has loops that cannot be contracted to a point. On the other hand, the special unitary group \(SU(n)\) is simply connected, because it does not contain any loops that cannot be contracted to a point.

    Another example is the special orthogonal group \(SO(3)\) in three dimensional space: For \(n \geq 3\), \(SO(n)\) is not simply connected. However, for \(n = 2\), \(SO(2)\) is simply connected. If we think about \(SU(2)\), which is the double cover of \(SO(3)\), we can see that it is simply connected: the inehrited loops from \(SO(3)\) can be contracted to a point in \(SU(2)\) because of the double cover structure. In fact, \(SU(2)\) is topologically equivalent to the 3-sphere \(S^3\), which is simply connected. Thus spin-\(\tfrac{1}{2}\) objects, which are described by representations of \(SU(2)\), can be thought of as living on a simply connected space, while the rotations in three-dimensional space, described by \(SO(3)\), have a non-simply connected structure due to the presence of loops that cannot be contracted to a point.
\end{example}

[Talking about foundamental groups, we are implying that we are working with a connected manifold, since the fundamental group is defined for connected spaces. If the space is not connected, we can consider the fundamental group of each connected component separately.]

\subsection{Lie Algebra from Tangent Vectors}

\begin{definition}[Tangent vectors]
    Let \((\mathcal{M}, \{u_{\alpha}\})\) and a point \(p \in \mathcal{M}\) in the manifold. Let us imagine a smooth curve \(\gamma : (a, b) \to \mathcal{M}\) such that \(\gamma(t=0) = p\). We can define the tangent vector to the curve at the point \(p\), \(\dot{\gamma}_p\) as as a differential operator acting on smooth functions \(\phi \in \mathbb{C}^{\infty}(\mathcal{M})\) as follows:
    \[
        \dot{\gamma}_p(\phi) = \left. \frac{\mathrm{d}}{\mathrm{d} t} (\phi \circ \gamma) (t)\right|_{t=t_0}.
    \]
\end{definition}

A set of tangent vectors at a point \(p\) spans a \(m\) dimensional vector space (\(m = \mathrm{dim} \mathcal{M}\)), called the tangent space at \(p\), denoted by \(T_p \mathcal{M}\). The dimension of the tangent space is thus equal to the dimension of the manifold.

Let us now consider a chart \(p \in U\) with coordinates \(x \equiv x^i : U \to \mathbb{R}^m\). We can define a set of tangent vectors at \(p\) as follows:
\[
    \dot{\gamma}_p^i(\phi) = \left. \frac{\mathrm{d}}{\mathrm{d} t} (\phi \circ x^{-1} \circ x \circ \gamma) (t)\right|_{t=t_0} = \left. \frac{\mathrm{d}}{\mathrm{d} t} (\phi \circ x^{-1}) (x(\gamma(t)))\right|_{t=t_0}.
\]
Applying the chain rule, we can write this as:
\[\TODO{correct the expression below, it is not correct}
    \dot{\gamma}_p^i(\phi) = \left. \frac{\mathrm{d} x^j}{\mathrm{d} t} \frac{\partial}{\partial x^j} (\phi \circ x^{-1}) (x(\gamma(t)))\right|_{t=t_0} = \left. \frac{\mathrm{d} x^j}{\mathrm{d} t}\right|_{t=t_0} \left. \frac{\partial}{\partial x^j} (\phi \circ x^{-1}) (x)\right|_{x=x(p)}.
\]
Thus we find that
\[
    \dot{\gamma}_p(\phi) = \sum_{i} \gamma^i(t_0) \frac{\partial}{\partial x^i} (\phi \circ x^{-1}) (x(p)) = \sum_{i} \gamma^i(t_0) \frac{\partial}{\partial x^i} (\phi(p)).
\]
In a local coorfinate patch thus the tangent vector can be expressed as a linear combination of the partial derivatives with respect to the coordinates, evaluated at the point \(p\). The coefficients of this linear combination are given by the components of the curve \(\gamma\) at \(t=t_0\). Thus the tangent space has dimension equal to the number of coordinates, which is equal to the dimension of the manifold.

\begin{corollary}
    If we have two manifolds \(\mathcal{M}_1\) and \(\mathcal{M}_2\) of dimensions \(m_1\) and \(m_2\) respectively, then the tangent space of their product is the direct sum of their tangent spaces:
    \[
        T_p(\mathcal{M}_1 \times \mathcal{M}_2) \cong T_p\mathcal{M}_1 \oplus T_p\mathcal{M}_2 = T_p \mathcal{M},
    \]
    where \(p = (p_1, p_2)\) is a point in the product manifold \(\mathcal{M}_1 \times \mathcal{M}_2\).
\end{corollary}

\begin{definition}[Tangent bundle]
    If one take the union of all the tangent spaces at all points of the manifold, we obtain the \textbf{tangent bundle} \(T\mathcal{M} = \bigcup_{p \in \mathcal{M}} T_p \mathcal{M}\). The tangent bundle is itself a manifold of dimension \(2m\), where \(m\) is the dimension of the original manifold \(\mathcal{M}\) (we have to keep the original maniforld in the definition of the tangent bundle, since the tangent spaces are defined at each point of the manifold):
    \[
        T\mathcal{M} = \{ (p,\mathbf{v}) | p \in \mathcal{M}, \mathbf{v} \in T_p \mathcal{M} \},
    \]
    and furthermore we have a natural projection map \(\pi : T\mathcal{M} \to \mathcal{M}\) defined by \(\pi(p,\mathbf{v}) = p\). The tangent bundle is a vector bundle, since each fiber (the tangent space at each point) is a vector space.
\end{definition}

\begin{definition}[Vector fields]
    A \textbf{vector field} on a manifold \(\mathcal{M}\) is a smooth section of the tangent bundle, that is, a smooth map \(V : \mathcal{M} \to T\mathcal{M}\) such that \(\pi \circ V = \mathrm{id}_{\mathcal{M}}\). In other words, for each point \(p \in \mathcal{M}\), the vector field assigns a tangent vector \(V(p) \in T_p \mathcal{M}\). It is also called a \textit{section} of the tangent bundle \(T\mathcal{M}\).
\end{definition}

Thus we associate to each point of the manifold a tangent space, and a vector field is a smooth assignment of a tangent vector to each point of the manifold. We are cutting the fibers of the bundle, and we are left with a single vector at each point of the manifold.

Since locally this vector field is a map from the manifold to the tangent space, if and only if the map is smooth, then the vector field is smooth and can be written as \(\Gamma(T\mathcal{M})\).





\[
    \dot{\gamma}_p(\phi) \in T_p \mathcal{M},
\]

\[
    \dot{\gamma}_p(\phi) = \frac{\mathrm{d}}{\mathrm{d} t} (\phi \circ \gamma) (t_0) = \sum_{i} \gamma^i(t_0) \frac{\partial}{\partial x^i} (\phi(p)),
\]
where \(\partial_i = \frac{\partial}{\partial x^i}\) is the basis for the tangent space at the point \(p\). Thus we can express the tangent vector as a linear combination of the basis vectors of the tangent space, with coefficients given by the components of the curve \(\gamma\) at \(t=t_0\). This shows that the tangent space at each point of the manifold is a vector space, and that the dimension of the tangent space is equal to the dimension of the manifold.

The vector field is a smooth assignment of a tangent vector to each point of the manifold
\[
    V : \mathcal{M} \to T\mathcal{M}, \quad p \mapsto (p, V_p) \in T_p \mathcal{M},
\]
and it can be thought of as a smooth map from the manifold to the tangent bundle. The vector field can be expressed in local coordinates as a linear combination of the basis vectors of the tangent space, with coefficients that are smooth functions on the manifold:
\[
    x^i (p) \in \mathbb{R}^n, \quad x^i (p) \mapsto (x^i (p), V^i(p)) \in T_p \mathcal{M} \equiv \mathbb{R}^n \times \mathbb{R}^n,
\]
This allows us to perform calculus on the manifold using the vector fields, and to define concepts such as divergence, curl, and gradient in a coordinate-independent way. If this mapping is smooth for every choice of chart, then the vector field is smooth and can be written as \(\Gamma(T\mathcal{M})\), and it forms a Lie algebra under the Lie bracket operation defined by the commutator of vector fields. Thus we managed to generalize the notion of vector fields from Euclidean space to arbitrary manifolds, and we can use this to define the Lie algebra of a Lie group as the tangent space at the identity element of the group.

We can also apply a vector field \(V \in \Gamma(T\mathcal{M})\) on a function \(\phi \in \mathbb{C}^{\infty}(\mathcal{M})\) to obtain a new function \(V(\phi) : \mathcal{M} \to \mathbb{R}\) defined by
\[
    p \mapsto V(\phi)(p) := V_p(\phi) \in \mathbb{R}, \quad V(\phi) \in \mathbb{C}^{\infty}(\mathcal{M}),
\]
where \(V_p\) is the tangent vector assigned to the point \(p\) by the vector field \(V\). This allows us to define the action of a vector field on functions, and to perform calculus on the manifold using the vector fields.

If we now take two vector fields \(V, W \in \Gamma(T\mathcal{M})\), we can apply the previous construction to the function \(V(\phi)\) to obtain a new function \(W(V(\phi)) : \mathcal{M} \to \mathbb{R}\) defined by
\[
    W(V(\phi))(p) := W_p(V(\phi)) \in \mathbb{R}, \quad W(V(\phi)) \in \mathbb{C}^{\infty}(\mathcal{M}),
\]
where \(W_p\) is the tangent vector assigned to the point \(p\) by the vector field \(W\). This can be computed as follows:
\[
    W_p(V(\phi)) = \sum_{i} W_p^i \partial_i (V(\phi))(p) = \sum_{ij} W_p^i \partial_i (V^j(p) \partial_j (\phi \circ x^{-1}))(x(p)),
\]
where we have used the local coordinate expression for the vector fields; we can continue to compute this expression by applying the product rule for differentiation:
\[
    W_p(V(\phi)) = \sum_{ij} W_p^i \partial_i (V^j(p)) \partial_j (\phi \circ x^{-1})(x(p)) + \sum_{ij} W_p^i V_p^j \partial_i \partial_j (\phi \circ x^{-1})(x(p)).
\]
Note that a second derivative term has appeared in the expression, which we do not now how to interpret in terms of the original vector fields \(V\) and \(W\): this new term is not a vector anymore, it is not an object behaving as a vector field \(X(\phi)\) with \(X \in \Gamma(T\mathcal{M})\). We cannot interpret this result as another vector field operating on the same function \(\phi\), since the second derivative term does not have the same structure as the original vector fields.

\begin{definition}[Lie bracket of vector fields]
    The Lie bracket of two vector fields (or commutator) \(V, W \in \Gamma(T\mathcal{M})\) is defined as the \uline{new vector field} \([V, W] \in \Gamma(T\mathcal{M})\) such that
    \[
        [V, W](\phi) := V(W(\phi)) - W(V(\phi)),
    \]
    for all smooth functions \(\phi \in \mathbb{C}^{\infty}(\mathcal{M})\). This new vector field is defined by the difference of the two second derivative terms that appeared in the previous expression, and it can be shown that it satisfies the properties of a Lie bracket, as we will see.
\end{definition}

With this new object defined, we have an operation that takes two vector fields and produces a new vector field, with in charts can be written as
\[
    [V, W]^j = \sum_{i} \left( V^i \partial_i W^j - W^i \partial_i V^j \right) \partial_j.
\]

\begin{proposition}
    For any three vector fields \(X, Y, Z \in \Gamma(T\mathcal{M})\), the Lie bracket satisfies the following properties:
    \begin{itemize}
        \item \textbf{Antisymmetry}: \([X, Y] = -[Y, X]\).
        \item \textbf{Bilinearity}: \([aX + bY, Z] = a[X, Z] + b[Y, Z]\) and \([Z, aX + bY] = a[Z, X] + b[Z, Y]\) for all scalars \(a, b \in \mathbb{R}\).
        \item \textbf{Jacobi identity}: \([X, [Y, Z]] + [Y, [Z, X]] + [Z, [X, Y]] = 0\).
    \end{itemize}
\end{proposition}

\begin{corollary}
    With the previously defined Lie bracket, the set of vector fields \(\Gamma(T\mathcal{M})\) can be endowed with the structure of a \textit{non-associative} algebra over \(\mathbb{R}\). We said non-associative because the Lie bracket does not satisfy the associative property, but it satisfies the Jacobi identity instead:
    \[
        [[X,Y],Z] \neq [X,[Y,Z]].
    \]
    The Jacobi identity is a weaker condition than associativity, and it ensures that the Lie bracket defines a consistent algebraic structure on the space of vector fields.
\end{corollary}

It is now time to connect this construction to the Lie algebra of a Lie group.

\begin{definition}
    Given a smooth map \(f : \mathcal{M} \to \mathcal{N}\) between two manifolds, we can define an induced map (which will be identified with a \textit{vector space homomorphism} for every point in the manifold) on the tangent spaces, called the \textbf{pushforward} of \(f\), denoted by \((f_*)_p : T_p \mathcal{M} \to T_{f(p)} \mathcal{N}\). The pushforward is defined by
    \[
        (f_*)_p (\dot{\gamma}_p) = \dot{(f \circ \gamma)}_{f(p)} \quad \forall p \in \mathcal{M},
    \]
    then let \(h \in \mathbb{C}^{\infty}(\mathcal{N})\), and \(v = \dot{\gamma}_p \in T_p\mathcal{M}\), then we have
    \[
        \implies \left[(f_*)_p(v)\right](h) = v(h \circ f).
    \]
    In coordinates \(x\) for \(\mathcal{M}\) and \(y\) for \(\mathcal{N}\), we can write the pushforward as
    \[
        [(f_*)_p]^i_k = \partial_k (y^i \circ f \circ x^{-1})(x(p)) = (\mathrm{d} f)^i_k (x(p)),
    \]
    where \(\mathrm{d} f\) is the Jacobian matrix of the map \(f\) at the point \(p\): \(f^i(x^k)\). This is just the jacobian matrix of a vector-valued function, which is a linear map that takes a tangent vector at \(p\) and maps it to a tangent vector at \(f(p)\) by applying the derivative of the function \(f\) to the components of the tangent vector.
\end{definition}

This concept is known in the field of general relativity as the \textbf{Lie dragging} of a vector field along a curve, and it is used to describe how vector fields change as they are transported along curves in a manifold. The pushforward allows us to compare vector fields at different points in the manifold, and to understand how they transform under smooth maps between manifolds.

If the function \(f\) is a diffeomorphism (smooth and invertible), then the pushforward is an isomorphism between the tangent spaces at \(p\) and \(f(p)\), and it allows us to transfer the structure of the tangent space from one point to another. This is particularly useful when we want to define the Lie algebra of a Lie group, as we will see successively.

\begin{proposition}
    If \(f : \mathcal{M} \to \mathcal{N}\) is a diffeomorphism, then the pushforward \((f_*)_p\) is an isomorphism between vector fields on \(\mathcal{M}\) and vector fields on \(\mathcal{N}\):
    \[
        f_* : \Gamma(T\mathcal{M}) \to \Gamma(T\mathcal{N}), \quad V \mapsto f_*(V),
    \]
    \[
        (f_*(V))_{f(p)} = (f_*)_p(V_p) \in T_{f(p)} \mathcal{N},
    \]
    and since consequently the jacobian matrix of \(f\) is invertible, we can also define the inverse pushforward \((f^{-1}_*)_{f(p)} : T_{f(p)} \mathcal{N} \to T_p \mathcal{M}\) as the inverse of the pushforward, and it is also an isomorphism between vector fields on \(\mathcal{N}\) and vector fields on \(\mathcal{M}\).
\end{proposition}

We can now imagine to concatenate functions; let \(f_1 \in \mathbb{C}^{\infty}(\mathcal{M}, \mathcal{N})\) and \(f_2 \in \mathbb{C}^{\infty}(\mathcal{N}, \mathcal{Q})\) be two diffeomorphisms, then their composition \(f_2 \circ f_1 : \mathcal{M} \to \mathcal{Q}\) is also a diffeomorphism, and we can define the pushforward of the composition as follows:
\[
    \forall p \in \mathcal{M}, \quad [(f_2 \circ f_1)_*]_p = [(f_2)_*]_{f_1(p)} \circ [(f_1)_*]_p : T_p \mathcal{M} \to T_{f_2(f_1(p))} \mathcal{Q},
\]
which can be read as the pushforward of the composition is the composition of the pushforwards. This property is known as the \textbf{functoriality} of the pushforward, and it ensures that the pushforward behaves well under composition of functions, which is a fundamental property for any mathematical construction that aims to capture the structure of smooth maps between manifolds.

Now we are ready to define the Lie algebra of a Lie group as the tangent space at the identity element of the group, and to show how the Lie bracket of vector fields induces a Lie bracket on this tangent space, which will give us the structure of a Lie algebra on the tangent space at the identity element. This is a crucial step in understanding the relationship between Lie groups and their associated Lie algebras, and it will allow us to study the properties of Lie groups using the tools of Lie algebras.

\begin{definition}
    Let \(G\) be a Lie group, and let \(e\) be the identity element of \(G\). Then any \(a \in G\) defines a ``\textit{left-translation}'' which maps any element \(g \in G\) to \(ag\):
    \[
        L_a : G \to G, \quad g \mapsto L_a g = ag.
    \]
    This left-translation is a diffeomorphism, since it is smooth and has a smooth inverse given by \((L_{a})^{-1} = L_{a^{-1}}\). Left translation are natural diffeomorphisms of the Lie group, and they allow us to transfer the structure of the tangent space at the identity element to any other point in the group. In particular, we can use left translations to define a vector field on the Lie group by translating a tangent vector at the identity element to every other point in the group.
\end{definition}

\begin{definition}
    Let us consider a class of vector fields \(V \in \Gamma(TG)\) on a Lie group \(G\) that are invariant under left translations; they are called \textbf{left-invariant vector fields}, and they satisfy the following property:
    \[
        \forall a \in G, \quad (L_a)_* V = V,
    \]
    which means that the vector field is unchanged under left translations. In other words, for any point \(g \in G\), the value of the vector field
    \[
        \forall a,g \in G, \quad [(L_a)_* V]_ag = [(L_a)_*]_g (V_g) = V_{ag},
    \]
    which reads as the value of the vector field at the point \(ag\) is equal to the pushforward of the vector field at the point \(g\) under the left translation by \(a\). This property ensures that the vector field is consistent with the group structure of \(G\).
\end{definition}

\begin{lemma}
    Given a tangent vector \(v \in T_e G\) (any point would have been fine, so we chose the identity element for convenience), there exists a unique left-invariant vector field \(V^{(v)} \in \Gamma(TG)\) by translating \(v\) to every point in the group using left translations, which satisfies the initial conditions
    \[
        V^{(v)}_e = v.
    \]
    Since it is a left-invariant vector field, for every point \(g \in G\) we can use left translation to define the value of the vector field at that point.
\end{lemma}

On a lie group we have a one to one correspondence between tangent vectors at a point (the identity element for convenience) and left-invariant vector fields, which allows us to transfer the structure of the tangent space at the identity element to the entire group
\[
    v \in T_e G, \quad v \mapsto V^{(v)} \in \Gamma(TG), \quad (V^{(v)})_g = (L_g)_* (v) \in T_g G.
\]

The other ingredient we need to define the Lie algebra of a Lie group is the closeness of the Lie bracket of vector fields, which we have already defined.

\begin{proposition}
    The set of left-invariant vector fields on a Lie group \(G\) is closed under the Lie bracket operation \([\cdot,\,\cdot]\), which means that if \(V, W \in \Gamma(TG)\) are left-invariant vector fields, then
    \[
        [V,W] \in \Gamma(TG) \text{ is left-invariant}.
    \]
\end{proposition}

In order to prove this proposition, we need to show another final lemma:

\begin{lemma}
    Let \(f : \mathcal{M} \to \mathcal{N}\) suriective and \(\phi \in \mathbb{C}^\infty(\mathcal{N})\) (so that \((\phi \circ f \in \mathbb{C}^\infty(\mathcal{M}))\)) and \(V \in \Gamma(T\mathcal{M})\) be a vector field on \(\mathcal{M}\) such that \((f_*)(V)\) is a well defined vector field on \(\mathcal{N}\), then we have
    \[
        V(\phi \circ f) = (f_*(V))(\phi) \circ f \iff V(\phi \circ f) = (f_*(V))_{f(p)} (\phi) = ((f_*)(V)(\phi) \circ f)(p).
    \]
    This means that ...
\end{lemma}

\begin{proof}
    Let \(\phi \in \mathbb{C}^\infty(G)\). Then, for any \(a,g \in G\) we have
    \[
        \begin{aligned}
            [(L_a)_* [V,W]]_{ag}(\phi) = \dots \\
            = [V,W]_g (\phi \circ L_a) = V_g (W(\phi \circ L_a)) - W_g (V(\phi \circ L_a)),
        \end{aligned}
    \]
    and since we have the previous lemma, we can write for a generic vector field \(U\) on \(G\)
    \[
        [(L_a)_* U]_ag = U_{ag} \implies U(\phi \circ L_a)(g) = [(L_a)_* U]_{ag}(\phi) = U_{ag}(\phi) = U(\phi)(L_a(g)) = (U(\phi) \circ L_a)(g).
    \]
    Thus we can write
    \[
        \begin{aligned}
            ((L_a)_* [V,W])_{ag}(\phi) = V_g (W(\phi) \circ L_a) - W_g (V(\phi) \circ L_a) \\
            = (((L_a)_*)_g V)_g (W(\phi)) - (((L_a)_*)_g W)_g (V(\phi))                    \\
            = [(L_a)_* V]_{ag} (W(\phi)) - [(L_a)_* W]_{ag} (V(\phi))                      \\
            = V_{ag} (W(\phi)) - W_{ag} (V(\phi)) = [V,W]_{ag}(\phi),
        \end{aligned}
    \]
    since both \(V\) and \(W\) are left-invariant vector fields, and thus we have shown that \([(L_a)_* [V,W]]_{ag}(\phi) = [V,W]_{ag}(\phi)\) for all \(a,g \in G\), which means that \([V,W]\) is left-invariant.
\end{proof}

Thus we have shown that the set of left-invariant vector fields on a Lie group \(G\) is closed under the Lie bracket operation, which means that we can define a Lie algebra structure on the tangent space at the identity element of the group by using the Lie bracket of left-invariant vector fields. This is a crucial step in understanding the relationship between Lie groups and their associated Lie algebras.

\begin{definition}
    Let \(G\) be a Lie group, and let \(e\) be the identity element of \(G\). The \textbf{Lie algebra} of \(G\) (or \textit{associated to} \(G\)), denoted by \(\mathfrak{g}\), is defined as the \(\mathbb{R}\)-vector space \(\mathfrak{g} = T_e G\) with the bilinear product given by the Lie bracket of left-invariant vector fields:
    \[
        [\cdot,\,\cdot]_\mathfrak{g} : \mathfrak{g} \times \mathfrak{g} \to \mathfrak{g}, \quad (v,w) \mapsto [V^{(v)},W^{(w)}]_\mathfrak{g} := [V^{(v)}, V^{(w)}]_e,
    \]
    where \((\mathfrak{g}, [\cdot,\cdot]_\mathfrak{g})\) satisfies antisymmetry, bilinearity, and the Jacobi identity, and thus it is abstractly defining a Lie algebra.
\end{definition}

\begin{definition}
    A vector space basis \(\{T_i\}_{i=1,\dots,n}\) with \(n = \text{dim}\mathfrak{g}\) of the Lie algebra \(\mathfrak{g}\) is called a \textbf{basis of generators} of the Lie algebra. Since the Lie brackets of elements of the Lie algebra can be expressed as linear combinations of the basis elements (since the algebra is closed under the Lie bracket), we can write
    \[
        [T_i,\,T_j] = \sum_{k=1}^n c_{ij}^k T_k,
    \]
    where \(c_{ij}^k\) are the \textbf{structure constants} of the Lie algebra \(\mathfrak{g}\) with the basis of generators \(\{T_i\}\).
\end{definition}

We can highlight some properties of the structure constants \(c_{ij}^k\) that follow from the properties of the Lie bracket:
\begin{itemize}
    \item They are antisymmetric in the first two (lower) indices: \(c_{ij}^k = -c_{ji}^k\), which follows from the antisymmetry of the Lie bracket.
    \item They satisfy the Jacobi identity: \(\sum_{l=1}^n (c_{ij}^l c_{lk}^m + c_{jk}^l c_{li}^m + c_{ki}^l c_{lj}^m) = 0\) for all \(i,j,k,m\), which follows from the Jacobi identity of the Lie bracket.
\end{itemize}
These comes directly from the properties of the Lie bracket and thus their application onto the generators of the Lie algebra. Furthermore an abstract Lie algebra can be defined by specifying a choice of basis \(\{T_i\}\) along with a set of structure constants \(c_{ij}^k\) that satisfy the above properties, without any reference to a specific Lie group.

As an easy example, we can consider the tensor \(\epsilon_{ij}^k\) and check that it respects the above properties: along with a good choice of basis, it can be used to define the structure constants of the Lie algebra \(\mathfrak{su}(2)\) associated to the Lie group \(SU(2)\).

We used real coefficients for the Lie algebra, but we could have used complex coefficients as well, and thus we would have defined a complex Lie algebra. The choice of real or complex coefficients depends on the context and the applications we have in mind, and it does not affect the general properties of the Lie algebra. The extension and generalization is straightforward, and it is just a matter of changing the field over which the vector space is defined.

\section{Lie Algebra of Matrix Groups}

If we consider matrices \(G \subseteq \mathrm{GL}(V)\), with \(V\) (\(n = \text{dim}V\)) a vector space over \(\mathbb{R}\) or \(\mathbb{C}\) (with a proper basis chosen), then \(\mathrm{GL}(V) \simeq \mathrm{GL}(n, \mathbb{R})\) or \(\mathrm{GL}(n, \mathbb{C})\), respectively (for simplicity let us stick to the real case, the complex case is just a straightforward generalization).

Taking than any element \(g \in G \subseteq \mathrm{GL}(n, \mathbb{R})\) can be thought of as an \(n \times n\) matrix
\[
    g = (g^{AB})_{AB}, \quad g^{AB} : G \to \mathbb{R}, \quad g \to g^{AB},
\]
for any fixed pair of indices \(A,B\). Thus we can think of the elements of the group as being labeled by a set of real numbers, which are the entries of the matrix.

\begin{definition}
    Let \(G \subset \mathrm{GL}(n, \mathbb{R})\) be a matrix group. Let \(v \in T_k G\) and \((U,x\equiv x^k)\), such that a test function \(\phi \in \mathbb{C}^{\infty}(G)\) with \(k = 1,\dots , \text{dim}G \neq n^2\) and a point \(h \in U\), then
    \[
        \begin{aligned}
            v(\phi) = \sum_{k} v^k \partial_k \phi \big|_{x(h)}                                                         \\
            = \sum_{k,A,B} v^k \frac{\partial g^{AB}}{\partial x^k} \bigg|_{x(h)} \frac{\partial \phi}{\partial g^{AB}} \\
            = \sum_{A,B} v^{AB} \frac{\partial \phi}{\partial g^{AB}},
        \end{aligned}
    \]
    where we can see how the summation over \(k\) defines a sort of \textit{tangent matrix}
    \[
        v^{AB} := \sum_{k} v^k \frac{\partial g^{AB}}{\partial x^k} \bigg|_{x(h)}, \quad \frac{\mathrm{d}}{\mathrm{d} t} (\gamma^{AB}(t)) \bigg|_{t=t_0} = v^{AB},
    \]
    which is the tangent vector to the curve \(\gamma^{AB}(t)\) in the space of matrices, and it can be thought of as an element of the Lie algebra associated to the matrix group \(G\). Thus we can identify the tangent space at the identity element of the group with a subspace of the space of matrices, and this subspace is precisely the Lie algebra of the group.
\end{definition}





\begin{definition}
    Let \(G \subset \mathrm{GL}(n, \mathbb{R}) \subset \mathbb{R}^{n\times n}\) be a matrix group, where the groups from left to right are included injectively. Then there exists a \(g \in G\) such that \(g \mapsto (g^{AB})\in \mathbb{R}^{n^2}\) is a smooth and injective map, and thus we can identify the elements of the group with points in \(\mathbb{R}^{n^2}\).

    Let now \(v \in T_h G\) then \((g^{AB})_*(v) \in T_{h^{AB}}\mathbb{R}^{n \times n} \simeq \mathbb{R}^{n^2}\) is injective. We will denote \((g^{AB})_*(v) = v^{AB}\) is the tangent matrix:
    \[
        \{v^{AB}\} \subseteq T_{h^{AB}} \mathbb{R}^{n \times n} \simeq \mathbb{R}^{n^2}.
    \]
\end{definition}

For \(v = \dot{\gamma}(t_0)\) with \(\gamma : \mathbb{R} \to G\) a curve with \(\gamma(t_0) = h\), the tangent vector \(v\) corresponds to a tangent matrix \(v^{AB}\) in \(\mathbb{R}^{n^2}\). We have \(\forall \phi \in \mathbb{C}^{\infty}(G)\),
\[
    [(g^{AB})_*(v)](\phi) = \frac{\mathrm{d}}{\mathrm{d} t} (\phi \circ g^{AB} \circ \gamma) (t_0) = \frac{\mathrm{d}}{\mathrm{d} t} (g^{AB} \circ \gamma) (t_0) \frac{\partial \phi}{\partial g^{AB}},
\]
which implies
\[
    v^{AB} = \frac{\mathrm{d}}{\mathrm{d} t} (g^{AB} \circ \gamma) (t_0).
\]
[...]

\begin{example}[General Linear group]
    Let us consider \(G = \mathrm{GL}(n, \mathbb{R}) = \{M \in \mathbb{R}^{n\times n} : \det(M) \neq 0\}\), we are interested in finding the Lie algebra of the grop as the tangent space at the identity element \(e = \mathbb{I}_n\) of the group. We can identify the elements of the group with points in \(\mathbb{R}^{n^2}\) by mapping each matrix \(M\) to its entries \(M^{AB}\). Locally, around the identity element, the tangent space can be identified with the space of all \(n \times n\) matrices:
    \[
        T_e G \simeq \mathfrak{g} = T_e \mathbb{R}^{n \times n} \simeq \mathbb{R}^{n^2},
    \]
    and since the dimension of the algebra, the group, the manifold etc. coincides, we know it to be of dimension \(n^2\). When passing to the complex counterpart, we have \(G = \mathrm{GL}(n, \mathbb{C})\) and the Lie algebra is \(\mathfrak{g} = T_e G \simeq \mathbb{C}^{n^2}\), which is of dimension \(2n^2\) as a real vector space.
\end{example}

\begin{example}[Special Linear group]
    If we consider the special linear group \(G = \mathrm{SL}(n, \mathbb{R}) = \{M \in \mathrm{GL}(n, \mathbb{R}) : \det(M) = 1\}\), we can identify the elements of the group with points in \(\mathbb{R}^{n^2}\) by mapping each matrix \(M\) to its entries \(M^{AB}\). Let \(M(t) : \mathbb{R} \to G\) be a curve in the group with \(M(0) = e = \mathbb{I}_n\). Then the tangent vector at the identity is given by
    \[
        v = \frac{\mathrm{d}}{\mathrm{d} t} M(t) \bigg|_{t=0}.
    \]
    Since \(M(t)\) is a matrix in \(\mathrm{SL}(n, \mathbb{R})\), we have \(\det(M(t)) = 1\) for all \(t\). Taking the derivative of this condition with respect to \(t\) and evaluating at \(t = 0\) gives
    \[
        \frac{\mathrm{d}}{\mathrm{d} t} \det(M(t)) \bigg|_{t=0} = 0.
    \]
    Using the formula for the derivative of the determinant
    \[
        \frac{\mathrm{d}}{\mathrm{d} t} \det(M(t)) \bigg|_{t=0} = \Tr\left[\det(M(t)) (M^{-1}(t))^T \dot{M}(t)\right]_{t=0} = \Tr(\dot{M}(0)) = \Tr(v),
    \]
    this implies
    \[
        \mathrm{tr}(v) = 0,
    \]
    which means that the tangent vector \(v\) (derivative of an element) at the identity element must be a traceless matrix. Therefore, the Lie algebra of the special linear group is the set of all traceless \(n \times n\) matrices, which can be identified with a subspace of \(\mathbb{R}^{n^2}\):
    \[
        \mathfrak{sl}(n, \mathbb{R}) = T_e G \simeq \{X \in \mathbb{R}^{n \times n} : \mathrm{tr}(X) = 0\}.
    \]
    The dimension of this Lie algebra is \(n^2 - 1\), since the condition of being traceless reduces the number of independent parameters by one (but is the same as the original group, since there we had the condition on the determinant). In the complex case, we have \(G = \mathrm{SL}(n, \mathbb{C})\) and the Lie algebra is \(\mathfrak{sl}(n, \mathbb{C}) = T_e G \simeq \{X \in \mathbb{C}^{n \times n} : \mathrm{tr}(X) = 0\}\), which is of dimension \(2(n^2 - 1)\) as a real vector space.
\end{example}

\begin{example}[Orthogonal group]
    Consider the orthogonal group \(G = \mathrm{O}(n, \mathbb{R}) = \{M \in \mathbb{R}^{n \times n} : M^T M = \mathbb{I}_n\}\). To find its Lie algebra, we look at the tangent space at the identity element \(e = \mathbb{I}_n\). A curve \(M(t) : \mathbb{R} \to G\) with \(M(0) = \mathbb{I}_n\) satisfies \(M(t)^T M(t) = \mathbb{I}_n\) for all \(t\). Differentiating this condition with respect to \(t\) and evaluating at \(t = 0\) gives
    \[
        \frac{\mathrm{d}}{\mathrm{d} t} (M(t)^T M(t)) \bigg|_{t=0} = 0.
    \]
    Using the product rule,
    \[
        \dot{M}(0)^T \mathbb{I}_n + \mathbb{I}_n^T \dot{M}(0) = 0,
    \]
    which simplifies to
    \[
        \dot{M}(0)^T + \dot{M}(0) = 0.
    \]
    This means that the tangent vector \(v = \dot{M}(0)\) must be skew-symmetric. Therefore, the Lie algebra of the orthogonal group is the set of all skew-symmetric \(n \times n\) matrices:
    \[
        \mathfrak{o}(n, \mathbb{R}) = T_e G \simeq \{X \in \mathbb{R}^{n \times n} : X^T + X = 0\}.
    \]
    The dimension of this Lie algebra is \(n(n-1)/2\), as there are \(n(n-1)/2\) independent parameters for a skew-symmetric matrix.
\end{example}

\begin{example}[Special Orthogonal group]
    For the special orthogonal group \(G = \mathrm{SO}(n, \mathbb{R}) = \{M \in \mathrm{O}(n, \mathbb{R}) : \det(M) = 1\}\), the Lie algebra is the same as that of the orthogonal group, since the condition of having determinant equal to one does not impose any additional constraints on the tangent vectors at the identity element. Therefore, we have
    \[
        \mathfrak{so}(n, \mathbb{R}) = T_e G \simeq \{X \in \mathbb{R}^{n \times n} : X^T + X = 0\},
    \]
    which is of dimension \(n(n-1)/2\).

    We can study this group as the intersection of the special linear group and the orthogonal group, since it is defined by the conditions of being both orthogonal and having determinant equal to one. Thus we can write
    \[
        \mathrm{SO}(n, \mathbb{R}) = \mathrm{SL}(n, \mathbb{R}) \cap \mathrm{O}(n, \mathbb{R}),
    \]
    and consequently we can write the Lie algebra as the intersection of the Lie algebras of the special linear group and the orthogonal group:
    \[
        \mathfrak{so}(n, \mathbb{R}) = \mathfrak{sl}(n, \mathbb{R}) \cap \mathfrak{o}(n, \mathbb{R}) = \{X \in \mathbb{R}^{n \times n} : X^T + X = 0, \mathrm{tr}(X) = 0\}.
    \]
    But since the condition of being skew-symmetric already implies that the trace is zero, we can simplify this to
    \[
        \mathfrak{so}(n, \mathbb{R}) = \{X \in \mathbb{R}^{n \times n} : X^T + X = 0\} = \mathfrak{o}(n, \mathbb{R}),
    \]
    which is the same as the Lie algebra of the orthogonal group.
\end{example}

\begin{example}[Unitary group]
    For the unitary group \(G = \mathrm{U}(n, \mathbb{C}) = \{U \in \mathbb{C}^{n \times n} : U^\dagger U = \mathbb{I}_n\}\), we can find its Lie algebra by looking at the tangent space at the identity element \(e = \mathbb{I}_n\). A curve \(U(t) : \mathbb{R} \to G\) with \(U(0) = \mathbb{I}_n\) satisfies \(U(t)^\dagger U(t) = \mathbb{I}_n\) for all \(t\). Differentiating this condition with respect to \(t\) and evaluating at \(t = 0\) gives
    \[
        \frac{\mathrm{d}}{\mathrm{d} t} (U(t)^\dagger U(t)) \bigg|_{t=0} = 0.
    \]
    Using the product rule,
    \[
        \dot{U}(0)^\dagger \mathbb{I}_n + \mathbb{I}_n^\dagger \dot{U}(0) = 0,
    \]
    which simplifies to
    \[
        \dot{U}(0)^\dagger + \dot{U}(0) = 0.
    \]
    This means that the tangent vector \(v = \dot{U}(0)\) must be anti-Hermitian. Therefore, the Lie algebra of the unitary group is the set of all anti-Hermitian \(n \times n\) matrices:
    \[
        \mathfrak{u}(n, \mathbb{C}) = T_e G \simeq \{X \in \mathbb{C}^{n \times n} : X^\dagger + X = 0\}.
    \]
    The real dimension of this Lie algebra is \(n^2\), as there are \(n^2\) independent parameters for an anti-Hermitian matrix: the upper half of the matrix has \(\frac{n(n-1)}{2}\) independent parameters, and the diagonal elements (all purely imaginary since they have to be changed in sign when daggered) have \(n\) independent parameters, for a total of \(2 \frac{n(n-1)}{2} + n = n^2\) independent parameters.
\end{example}

\begin{example}[Special Unitary group]
    For the special unitary group \(G = \mathrm{SU}(n, \mathbb{C}) = \{U \in \mathrm{U}(n, \mathbb{C}) : \det(U) = 1\}\), the Lie algebra is the set of all anti-Hermitian \(n \times n\) matrices with trace zero:
    \[
        \mathfrak{su}(n, \mathbb{C}) = T_e G \simeq \{X \in \mathbb{C}^{n \times n} : X^\dagger + X = 0, \mathrm{tr}(X) = 0\}.
    \]
    The real dimension of this Lie algebra is \(n^2 - 1\), as the condition of being traceless reduces the number of independent parameters by one.
\end{example}

We are computing the \textbf{real dimension} of the lie algebra even in the complex case, since we are considering the Lie algebra as a real vector space. If we were to consider it as a complex vector space, we would have to consider that when the coefficients are complex, the dimension would be halved, since each complex parameter corresponds to two real parameters (the real and imaginary parts). Thus, when we have an odd number of independent parameters, the complex dimension would be a half-integer, which is not possible for a vector space. Therefore, we consider the Lie algebra as a real vector space, even in the complex case, to ensure that the dimension is always an integer.

We want to impose only \textbf{holomorphic functions} on the complex Lie groups, since we want to preserve the complex structure of the group and its Lie algebra. If we were to allow non-holomorphic functions, we would lose the complex structure and we would not be able to define a complex Lie algebra. The holomorphic functions (for example the determinant consition, which is just a polynomial equation in complex variables) ensure that the Lie algebra is closed under the Lie bracket and that it has a well-defined complex structure, which is essential for many applications in mathematics and physics.

\begin{lemma}
    Let \(G\) be a matrix Lie group, and let \(v \in T_h G\) be a tangent vector in \(h\) and consider \(a \in G\); then the pushforward of \(v\) under the left translation by \(a\) is given by
    \[
        [[(L_a)_*]_h(v)]^{AB} = \sum_{C=1}^n a^{AC} v^{CB} = (a \cdot v)^{AB},
    \]
    where the last equality makes sense only in a matrix group, since we are using the matrix multiplication to define the action of the group on the tangent vector. This shows how the pushforward of a tangent vector under left translation can be expressed in terms of the matrix multiplication of the group element and the tangent vector, which is a key property that allows us to understand how the Lie algebra is related to the structure of the Lie group.
\end{lemma}
\begin{proof}
    We can use a test function \(\phi \in \mathbb{C}^\infty(G)\) to compute the pushforward of \(v\) under left translation by \(a\):
    \[
        [(L_a)_*]_h(v)(\phi) = v(\phi \circ L_a) = \sum_{k} v^k \partial_k (\phi \circ L_a) \big|_{x(h)} \equiv v(\tilde{\phi}),
    \]
    where we defined \(\tilde{\phi} := \phi (ag)\). Now we can express the function \(\tilde{\phi}\) in terms of the matrix entries \(g^{AB}\) if we compute
    \[
        [[(L_a)_*]_h(v)]^{AB} = v^{AB} \frac{\partial \phi}{\partial g^{PQ}} \bigg|_{g = ah} \frac{\partial}{\partial g^{AB}} (ag)^{PQ} \bigg|_{g = h},
    \]
    which can be rewritten as
    \[
        = v^{AB} \frac{\partial \phi}{\partial g^{PQ}} \bigg|_{g = ah} a^{PR} \delta^{RA}\delta^{QB} = a^{PR} v^{RQ} \frac{\partial \phi}{\partial g^{PQ}} = (a \cdot v)^{AB} \frac{\partial \phi}{\partial g^{AB}},
    \]
    [...]
\end{proof}

\begin{proposition}
    Let \(G\) be the matrix group and \(v,w \in T_e G \implies \mathfrak{g}\), then we can associate to \(v\) and \(w\) the left-invariant vector fields \(V^{(v)}\) and \(W^{(w)}\), and we can compute the Lie bracket of these vector fields at the identity element:
    \[
        [V^{(v)}, W^{(w)}]_e = [v,w]_\mathfrak{g},
    \]
    where the Lie bracket on the right-hand side is defined by the commutator of the corresponding tangent matrices:
    \[
        [v,w]_\mathfrak{g} = (v w)^{AB} - (w v)^{AB} = v^{AC} w^{CB} - w^{AC} v^{CB},
    \]
    which is the standard Lie bracket for matrix Lie algebras and makes sense in the context of matrix groups, since the tangent vectors can be identified with matrices and the Lie bracket can be defined as the commutator of these matrices.
\end{proposition}

We are not going to prove this proposition, since it is a straightforward computation that follows from the definition of the Lie bracket of vector fields and the properties of the pushforward under left translations. The key point is that the left-invariant vector fields are defined by translating the tangent vectors at the identity element to every other point in the group using left translations, and thus their Lie bracket can be computed by using the properties of the pushforward under left translations and the definition of the Lie bracket of vector fields.

\section{Exponential Map}

We are now going to define the exponential map for a Lie group, which is a crucial tool for understanding the relationship between the Lie algebra and the Lie group. The exponential map allows us to relate elements of the Lie algebra (tangent vectors at the identity element) to elements of the Lie group, and it provides a way to construct curves in the Lie group from tangent vectors in the Lie algebra. It is the tool which let us take a finite transformation from an infinitesimal one, and thus it is fundamental for applications in physics, where we often deal with infinitesimal transformations and we want to understand how they relate to finite transformations.

\begin{definition}[Integral curve]
    Let \(V \in \Gamma(T \mathcal{M})\) be a vector field and let \(\gamma: (a,b) \to \mathcal{M}\) be an integral curve for \(V\), i.e.,a smooth curve whose tangent vector at each point \(\gamma(t)\) is equal to the value of the vector field at that point \(V_{\gamma(t)}\), i.e.
    \[
        \dot{\gamma}(t) = V_{\gamma(t)} \quad \forall t \in (a,b).
    \]
    If we use a chart \((U,x)\) around a point \(p \in \mathcal{M}\) such that \(\gamma(t_0) = p\) for some \(t_0 \in (a,b)\), then we can write the integral curve in coordinates as
    \[
        \dot{\gamma}^i(t) = \frac{\mathrm{d}}{\mathrm{d}t} \gamma^i(t) = V^i(\gamma(t)) \quad \forall t \in (a,b), \quad \text{(ODE)},
    \]
    which is a system of ordinary differential equations (ODEs) that can be solved to find the integral curve \(\gamma(t)\) given the vector field \(V\) and the initial condition \(\gamma(t_0) = p\).
\end{definition}

The existence and uniqueness of solutions to this ODE is guaranteed by the standard theory of ODEs, and thus we can always find an integral curve for any given vector field and initial condition.

\begin{definition}[One-parameter subgroup]
    Let a smooth curve \(\sigma : \mathbb{R} \to G\) in a Lie group \(G\) is called a \textbf{one-parameter subgroup} of \(G\) if it satisfies the group homomorphism property
    \[
        \sigma(s+t) = \sigma(s) \sigma(t) \quad \forall s,t \in \mathbb{R},
    \]
    with \(\sigma(0) = e\) is the identity element of the Lie group \(G\).
\end{definition}

This last definition means that a one-parameter subgroup is a curve in the Lie group that respects the group structure, i.e., the value of the curve at the sum of two parameters is equal to the product (group operation) of the values of the curve at each parameter. This property is crucial for understanding how one-parameter subgroups relate to the Lie algebra and how they can be used to construct finite transformations from infinitesimal ones.

We will adopt the notation in which
\[
    \sigma \equiv \{\sigma(t) \mid t \in \mathbb{R}\},
\]
as a subgroup of \(G\) generated by the one-parameter subgroup \(\sigma\). This notation emphasizes the fact that the one-parameter subgroup generates a subgroup of \(G\) that is isomorphic to \(\mathbb{R}\) under addition, and it allows us to think of the one-parameter subgroup as a continuous family of group elements parameterized by a real variable \(t\).

\begin{proposition}
    Let \(G\) be a Lie group and let \(V \in \Gamma(TG)\) be a left-invariant vector field on \(G\). Then the integral curve \(\sigma : \mathbb{R} \to G\) of \(V\) with initial condition \(\sigma(0) = e\) is a one-parameter subgroup of \(G\).
\end{proposition}

The idea here is that by taking the integral curve of a left-invariant vector field starting at the identity element, we obtain a curve that respects the group structure and thus defines a one-parameter subgroup. The left-invariance of the vector field ensures that the flow generated by the vector field is compatible with the group operation, which is essential for the resulting curve to be a one-parameter subgroup.

\begin{proposition}
    Let \(\sigma\) be a one-parameter subgroup of a Lie group \(G\). Then the tangent vector \(V_{\sigma(t)} = \frac{d}{dt}\sigma(t)\) to the curve \(\sigma(t)\) at any point \(t\) can be computed as the pushforward of the tangent vector at the identity element \(V_e\) under the left translation by \(\sigma(t)\):
    \[
        V_{\sigma(t)}  = [(L_{\sigma(t)})_*]_e (\dot{\sigma}(0)),
    \]
    which implies that \(V_{\sigma(t)}\) is a left-invariant vector field along the curve \(\sigma(t)\), and thus the one-parameter subgroup \(\sigma\) is generated by the left-invariant vector field \(V\) that corresponds to the tangent vector at the identity element.
\end{proposition}

Thus every one-parameter subgroup of a Lie group is generated by a left-invariant vector field, and conversely, every left-invariant vector field generates a one-parameter subgroup. This establishes a deep connection between the Lie algebra (the space of left-invariant vector fields) and the structure of the Lie group (the one-parameter subgroups), which is fundamental for understanding how the Lie algebra encodes information about the local structure of the Lie group and how it can be used to construct finite transformations from infinitesimal ones.

It is a one-to-one correspondence between the elements of the Lie algebra and the one-parameter subgroups of the Lie group, which is a key result in the theory of Lie groups and Lie algebras.

\begin{corollary}
    Suppose to have two different one-parameter subgroups \(\sigma(t)\) and \(\mu(t)\) of a Lie group \(G\) with \(\dot{\sigma}(0) = \dot{\mu}(0) = v \in T_e G\). Then \(\sigma(t) = \mu(t)\) for all \(t \in \mathbb{R}\).
\end{corollary}

This corollary can be read as saying that the one-parameter subgroup generated by a given tangent vector at the identity element is unique. This emphasizes the one-to-one correspondence between the elements of the Lie algebra and the one-parameter subgroups of the Lie group, and it ensures that each element of the Lie algebra generates a unique one-parameter subgroup.

\begin{definition}
    Let \(X \in \mathfrak{g}\) be an element of the Lie algebra of a Lie group \(G\). We denote as \(\sigma_X\) the unique one-parameter subgroup of \(G\) constructed as follows:
    \begin{enumerate}
        \item \(X \in \mathfrak{g} \iff V^{(X)} \in \Gamma(TG)\) is left-invariant vector field on \(G\) such that \(V^{(X)}_e = X\).
        \item \(\sigma_X : \mathbb{R} \to G\) is the unique integral curve of \(V^{(X)}\) with initial condition \(\sigma_X(0) = e\) (passing through the identity element) with \(\dot{\sigma}_X(0) = V_e^{(X)} = X\); it is a one-parameter subgroup of \(G\).
    \end{enumerate}
\end{definition}

Since along the curve \(\sigma_X(t)\) we have a possible mapping from the Lie algebra to the Lie group, we can define the exponential map as follows.

\begin{definition}[Exponential map]
    The exponential map
    \[
        \exp : \mathfrak{g} \to G, \quad X \mapsto \sigma_X(1),
    \]
    is called the exponential map of the Lie group \(G\). It maps an element \(X\) of the Lie algebra \(\mathfrak{g}\) to the value of the one-parameter subgroup \(\sigma_X\) at \(t = 1\), which is an element of the Lie group \(G\).
\end{definition}

\begin{proposition}
    If \(X \in \mathfrak{g}\), \(\lambda \in \mathbb{R}\) then we have
    \[
        \exp(\lambda X) = \sigma_X(\lambda).
    \]
    Since this is a one-parameter subgroup, we have
    \[
        \exp((\lambda + \mu) X) = \sigma_X(\lambda + \mu) = \sigma_X(\lambda) \sigma_X(\mu) = \exp(\lambda X) \exp(\mu X), \quad \forall \lambda, \mu \in \mathbb{R},
    \]
    and we can also derive straightforwardly that \(\exp(0) = e\) and \(\exp(-X) = \sigma_X(-1) = \sigma_X(1)^{-1} = \exp(X)^{-1}\). This shows that the exponential map has properties that are consistent with the group structure, and it allows us to relate the Lie algebra to the Lie group in a way that respects the algebraic structure of both.
\end{proposition}

This definition allows us to plug in the map any vector from the Lie algebra and obtain an element of the Lie group, which is a powerful tool for understanding how the Lie algebra encodes information about the local structure of the Lie group and how it can be used to construct finite transformations from infinitesimal ones.

\begin{proof}
    For \(\lambda \neq 0\) we can define a new curve \(\tilde{\sigma} : \mathbb{R} \to G\) as
    \[
        \tilde{\sigma}(t) := \sigma_X(\lambda t),
    \]
    which is still a one-parameter subgroup of \(G\) since it satisfies the group homomorphism property. If we now define another one-parameter subgroup from \(\sigma_{\lambda X}(0) = \lambda X \), then we have also
    \[
        \dot{\tilde{\sigma}}(0) = \frac{\mathrm{d}}{\mathrm{d} t} \sigma_X(\lambda t) \bigg|_{t=0} = \lambda \dot{\sigma}_X(0) = \lambda X,
    \]
    which means that \(\tilde{\sigma}\) and \(\sigma_{\lambda X}\) are two one-parameter subgroups of \(G\) with the same tangent vector at the identity element (initial conditions), and thus by the previous corollary they must coincide, i.e., \(\tilde{\sigma}(t) = \sigma_{\lambda X}(t)\) for all \(t \in \mathbb{R}\). If this is the case, then we can use the definition of the exponential map with respect to the one-parameter subgroup \(\sigma_{\lambda X}\) to write
    \[
        \exp(\lambda X) = \sigma_{\lambda X}(1) = \tilde{\sigma}(1) = \sigma_X(\lambda),
    \]
    which is the desired result.
\end{proof}

We are now going to express some basic and important properties of the exponential map.

\begin{proposition}
    The exponential map has the following properties:
    \begin{enumerate}
        \item The image of the exponential map is a neighborhood of the identity element \(e\) in \(G\), which means that
              \[
                  \exp(\mathfrak{g}) \subseteq G_0 = \{g \in G | \exists \gamma : [a,b] \to G \text{ continuous, with } \gamma(t_0) = e, \gamma(t_1) = g\},
              \]
              where \(G_0\) is the connected component of the identity element in \(G\). For example \(\mathrm{SO}(n) = [O(n)]_e\).
        \item The Lie algebra \(\mathfrak{g}\) is a \(\mathbb{R}\)-vector space \(\simeq \mathbb{R}^n\), then we can say that the exponential map is a smooth map from \(\mathbb{R}^n\) to \(G\), and thus it is a \textbf{local diffeomorphism} around the identity element \(e\). This means that there exists a neighborhood \(U\) of \(0 \in \mathfrak{g}\) such that \(\exp : U \to \exp(U)\) is a diffeomorphism onto its image, which is a neighborhood of the identity element in \(G\). With these restrictions the exponential map is thus injective.
        \item The exponential map \(\exp : \mathfrak{g} \to G\) is a smooth map, but generally not surjective, since there may be elements of the Lie group that cannot be reached by exponentiating elements of the Lie algebra. But if \(G_0\) is a compact (connected component) then it is surjective. For example, if \(G = \mathrm{SL}(2, \mathbb{R})\), which is not compact, then the exponential map is not surjective. On the other hand, if \(G = \mathrm{SU}(n)\), which is compact, then the exponential map is surjective. There may be cases of surjectivity even for non-compact groups, but this is not guaranteed in general (for example, the exponential map for \(\mathrm{GL}(n,\mathbb{C})\) and \(\mathrm{SO}(1,n)\) is surjective even if the group is not compact).
        \item Any \(g \in G_0\) can be written as
              \[
                  g = \exp(X_1) \exp(X_2) \cdots \exp(X_k),
              \]
              for some \(X_1, X_2, \ldots, X_k \in \mathfrak{g}\) and some \(k \in \mathbb{N}\). This means that any element of the connected component of the identity can be expressed as a finite product of exponentials of elements of the Lie algebra, which do not contradicts the previous point about surjectivity, since we are not claiming that every element of \(G\) can be written as a single exponential, but rather that it can be written as a finite product of exponentials.
        \item \textbf{Baker-Campbell-Hausdorff (BCH) formula}: For any \(X,Y \in \mathfrak{g}\), we have then
              \[
                  Z = X + Y + \frac{1}{2} [X,Y]_\mathfrak{g} + \frac{1}{12} [X,[X,Y]_\mathfrak{g}]_\mathfrak{g} - \frac{1}{12} [Y,[X,Y]_\mathfrak{g}]_\mathfrak{g} + \cdots,
              \]
              where the series continues with higher-order commutators. This formula let us connect ...
              Furtheore, if this series converges, we have
              \[
                  \exp(Z) = \exp(X) \exp(Y) = \exp\left(X + Y + \frac{1}{2} [X,Y]_\mathfrak{g} + \frac{1}{12} [X,[X,Y]_\mathfrak{g}]_\mathfrak{g} - \frac{1}{12} [Y,[X,Y]_\mathfrak{g}]_\mathfrak{g} + \cdots\right),
              \]
              has a solution for \(Z\) in terms of \(X\) and \(Y\). This formula is fundamental for understanding how the group operation in \(G\) can be expressed in terms of the Lie algebra \(\mathfrak{g}\), and it provides a way to compute the product of two exponentials in terms of the Lie bracket and higher-order commutators. If the group is compact or the map is surjective, then we can use the BCH formula to express any element of the group as an exponential of an element of the Lie algebra (since the series would always converge in this case). For the case in which \([X,Y]=0\) then the BCH formula simplifies to
              \[
                  \exp(X) \exp(Y) = \exp(X + Y),
              \]
              which is a special case that occurs when the Lie algebra is abelian (i.e., when all elements commute with each other).
        \item Let \(F : G_1 \to G_2\) be a smooth map between two Lie groups (Lie group homomorphism), then
              \[
                  \exp_2(F_*(X)) = F(\exp_1(X)), \quad \forall X \in \mathfrak{g}_1,
              \]
              where \(\exp_1\) and \(\exp_2\) are the exponential maps of \(G_1\) and \(G_2\) respectively, and \(F_*\) is the pushforward of \(F\). This property shows that the exponential map is compatible with Lie group homomorphisms, and it allows us to relate the exponential maps of different Lie groups through the pushforward of a homomorphism between them.
    \end{enumerate}
\end{proposition}

